\section{Groups and schemes}\label{section: groups and schemes}
This section serves as a recollection of concepts from algebraic geometry and representation theory which we freely use later in the paper.

\subsection{Schemes}

\subsubsection{Base rings.}
Let $k$ be a base ring for our geometry. The most relevant case is when $k$ is an algebraically closed field of $\chara k = 0$, for example $k=\C$. The reader is invited to focus on this case.

On the other hand, let $\k$ be an arbitrary ring of coefficients. Again, one may want to take $\k =\C$ for simplicity. Nevertheless, a majority of our results works over fields $\k$ of arbitrary characteristic. Even better, most of our constructions work over $\k = \Z$.  


\subsubsection{General conventions.}
Given a scheme $X$, we write $\O(X)$ for the ring of global sections of its structure sheaf $\O_X$.
We write $|X|$ for the set of closed points of $X$ -- we prefer to use this concise notation for closed points only for the sake of readability; this should not lead to any confusion.


In general, we write $\coprod$ for actual coproduct (disjoint union) of schemes. On the other hand, we write $\bigsqcup$ for locally closed stratifications.

\subsubsection{Closed subschemes and fibers.}
Given a closed subscheme $Z \hookrightarrow X$, we write $X - Z$ for the open complement. The ideal sheaf of $Z$ in $X$ is denoted by $\eI_{Z}$; it sits in the short exact sequence
\begin{equation*}
    0 \to \eI_{Z} \to \O_X \to \O_Z \to 0.
\end{equation*}
We further write $X_{Z, n} = {\Spec}_{\O_X}(\O_X / \eI^n_Z) \hookrightarrow X$ for the $n$-th infinitesimal thickening of $Z$ inside $X$. We denote $X^{\wedge}_Z := \lim_{n} X_{Z, n}$ the formal completion of $X$ along $Z$.


Given a map of schemes $X \to Y$ and a point $y \in Y$, we write $X_y$ for the fiber. Furthermore, we denote by $X^{\wedge}_{y}$ the formal completion of $X$ around $y$. By definition, $X^{\wedge}_y := X^{\wedge}_{X_y}$ is the completion of $X$ along the fiber $X_y$.

\subsubsection{Quotients.}
Let $G$ be a group scheme. Suppose $G \acts X$. We write $X/G$ for the stack quotient (taken in the fpqc topology). 

On the other hand, we denote by $X \sslash G$ the GIT quotient (in general this depends on further choices). We only consider GIT quotients for $X = \Spec R$ affine, when $X\sslash G$ is also affine with ring of functions $\O(X \sslash G) = \O(X/G) = \O(X)^G= R^G$.


\subsubsection{Ind-schemes.}
Consider the Yoneda embedding $\Sch_k \hookrightarrow \Sh(\Sch_k)$ into the category of sheaves (say of sets) on $\Sch_k$ (say in the étale topology). This embedding commutes with all limits, but not with colimits. The category $\Sch_k$ of $k$-schemes is in fact complete, but not cocomplete. In other words, colimits of schemes may only exists as objects of the bigger category $\Sh(\Sch_k)$. 

In particular, suppose $X$ is given as a filtered colimit $X = \colim_{i \in I} X_i$ of schemes $X_i$ along closed immersions. Such $X$ are called \textit{ind-schemes}, and span a full subcategory $\IndSch_k$. Any filtered system $(X_i, i \in I)$ as above together with a chosen isomorphism $X = \colim_{i \in I} X_i$ is called an \textit{ind-presentation} of $X$. Clearly, $\Sch_k \subseteq \IndSch_k$ via constant systems.


Many construction and properties of schemes and their morphisms extend to ind-schemes by checking on a ind-presentation. In particular, the notions of ind-proper and ind-closed immersion make sense. 

\subsubsection{Functions on ind-schemes.}
The global functions on an ind-scheme are given by the inverse limit of global functions on any ind-presentation
\begin{equation*}
    \O(X) := \Hom(X, \A^1) = \Hom(\colim_{i\in I} X_i, \A^1) = \lim_{i \in I}(X_i, \A^1) = \lim_{i \in I} \O(X_i).
\end{equation*}
Similarly for other representable functors in place of $\A^1$.


\begin{remark}
    Given a closed immersion of ind-schemes $Z \hookrightarrow X$, the induced map on global functions $\O(X) \to \O(Z)$ need not be surjective. For an illustrative example, let $X = \A^1_k$ over a field $k$ of characteristic zero and $Z = \coprod_{n \in \Z} \pt = \colim_{i \in \N} \coprod_{n \in [-i, i] } \pt$, with the embedding $Z \hookrightarrow X$ given by the discrete set of integer points in $\A^1(k)$. This is clearly a closed ind-subscheme, since the map at each finite stage is a closed immersions of schemes. However, the induced map on global functions $\O(\A^{1}) \to \prod_{n \in \Z} k$ is not surjective: not every function $\Z \to k$ may be interpolated by a polynomial.

    Given a closed immersion $Z \hookrightarrow X$ of ind-schemes such that $Z$ is a scheme, the map $\O(X) \to \O(Z)$ is surjective. 
\end{remark}






\subsection{Deformations to normal cones and affine blowups}\label{section: deformations to normal cones and affine blowups}


\subsubsection{Blowups.}\label{section: blowups}
Given a scheme $X$ with a closed subscheme $Z \hookrightarrow X$, we denote $\Bl_Z(X)$ the blowup of $X$ at $Z$. It has a natural map $\pi: \Bl_Z(X) \to X$. The fiber $E:=\pi^{-1}(Z) \subseteq \Bl_Z(X)$ is an effective Cartier divisor. The pair $(\Bl_Z(X), E)$ is the universal pair mapping to $X$ characterized by the property that $E \hookrightarrow \Bl_Z(X)$ is an effective Cartier divisor \cite[Tag 085U]{Sta}.

In particular $\Bl_X(X) = \emptyset$. In an ambient scheme $X$, suppose we are given closed subschemes $Z \hookrightarrow X_1, X_2 \hookrightarrow X$. Inside the ambient blowup $\Bl_Z(X)$, we then have
\begin{equation*}
    \Bl_{Z}(X_1) \cap \Bl_{Z}(X_2) = \Bl_{Z}(X_1 \cap X_2).
\end{equation*}

\subsubsection{Strict transforms.}\label{section: strict transforms}
See \cite[Tag 080C]{Sta} for a general discussion. Consider a closed subscheme $Z \hookrightarrow X$, and the associated blowup $\Bl_Z(X)$. Given any other closed subscheme $Y \hookrightarrow X$, we have its strict transform
\begin{equation*}
    Y^{\st} \hookrightarrow \Bl_Z(X).
\end{equation*}
which may be described as the scheme-theoretic closure of the preimage of $Y \cap (X - Z)$, namely
\begin{equation*}
    Y^{\st} = \overline{\pi^{-1}(Y \cap (X - Z))} \hookrightarrow \Bl_Z(X).
\end{equation*}
%
In terms of $Y$, the strict transform is given as
\begin{equation*}
    Y^{\st} = \Bl_{Y \cap Z} Y \hookrightarrow \Bl_{Z} X.
\end{equation*}
Here, the embedding comes from the functoriality of blowups. See \cite[Tag 080E]{Sta}.



\subsubsection{Deformations to normal cones.}
See \cite[\S 5]{Ful84}.
Let $Z \hookrightarrow X$ be a closed subscheme. Consider the embedding $Z \times \{0\} \hookrightarrow X \times \A^1$. We can consider the blowup $\Bl_{Z \times \{0\}}(X \times \A^1)$. The \textit{deformation of $X$ to the normal cone of $Z$} is then given by the open complement to the strict transform of $X \times \{0\}$, namely
\begin{equation*}
    N(Z, X) := \Bl_{Z \times \{0\}}(X \times \A^1) - (X \times \{0\})^{\st}.
\end{equation*}
%
It is an affine flat scheme over $\A^1$. Denoting by $C(Z, X)$ the normal cone of $Z \subseteq X$, meaning the relative spectrum of the sheaf of algebraic differentials $\Omega_{Z/X}$, the deformation $N(Z, X)$ fits in the pullback diagram
\begin{equation*}
    \begin{tikzcd}
      C(Z, X) \arrow[d] \arrow[r]&  N(Z,X) \arrow[d] & \arrow[l] \arrow[d] X \times (\A^1 - \{ 0 \}) \\
      \{0\} \arrow[r] & \A^1 & \arrow[l] \A^1 - \{ 0 \}. 
    \end{tikzcd}
\end{equation*}
In words, $N(Z, X)$ carries a structure map $N(Z,X) \to \A^1 = \Spec(k[h])$, whose fibers over $h \neq 0$ are isomorphic to $X$, while the fiber over $h=0$ is given by the normal cone $C(Z, X)$ of $Z$ inside $X$.


More explicitly, the structure sheaf $\O_X$ has a descending filtration $(\eI^j_{Z})_{j = 0}^{\infty}$ by powers of the ideal sheaf $\eI_Z$. The associated Rees construction
\begin{equation*}
 \O_X[\tfrac{f}{h^j} \mid f \in \eI^j_Z]    
\end{equation*}
turns it into a graded algebra, whose relative spectrum gives the above deformation:
\begin{equation*}
 N(Z, X) = \Spec_{\O_X}(\O_X[\tfrac{f}{h^j} \mid f \in \eI^j_Z]).  
\end{equation*}






\subsubsection{Basic properties.}

The deformation to the normal cone is functorial in the input morphism $Z \to X$.
It is further compatible with pullbacks: given $Z \hookrightarrow X$ and $Y \to X$, we have a natural identification $N(Z, X) \times_X Y = N(Z\times_X Y, Y)$.
If a finite group $W$ acts on $Z \hookrightarrow X$, it also acts on $N(Z, X)$ and we have $N(Z \sslash W, X\sslash W) = N(Z, X) \sslash W$.

    
Suppose we are given a finite set of generators of the ideal sheaf $\eI_Z = \langle f_1, \dots, f_m  \rangle$. Then 
\begin{equation*}\label{section: generators of the deformation}
    \O_X[\tfrac{f}{h^j} \mid f \in \eI^j_Z] = \O_X[\tfrac{f_i}{h} \mid i = 1, \dots, m]. 
\end{equation*}
In other words, to generate the structure sheaf of $N(Z, X)$, we only need to add finitely many ring generators
$\tfrac{f_i}{h}$ for $i = 1, \dots, m$.




\subsubsection{Affine blowups.}\label{subsubsection: affine blowups.}
More generally, suppose we have a family of schemes $X \to \A^1$ over the base $\A^1$. Let $\zeta \in \A^1$ be a point, and suppose $Z \hookrightarrow Y_{\zeta}$ is a closed subscheme of the fiber of $X$ at $\zeta$.
Then the \textit{affine blowup of $X$ at $Z$ with respect to $X \to \A^1$}, denoted $\Bl^{\circ}_Z(X)$ is the open complement in $\Bl_{Z}(X)$ of the strict transform of $X_{\zeta}$, namely
\begin{equation*}
    \Bl^{\circ}_Z(X) := \Bl_Z(X) - X^{\st}_{\zeta}.
\end{equation*}
We have a natural affine morphism $\Bl^{\circ}_Z(X) \to X$.
%
By design, the deformation to the normal cone is a special case: $N(Z, X) = \Bl^{\circ}_{Z \times \{0\}}(X \times \A^1)$.

\subsubsection{Strict transforms inside affine blowups.}\label{section: embedding of strict transform into affine blowup}
Assume that we have another closed subscheme $Y \hookrightarrow X$ such that $Z = Y \cap X_{\zeta}$. Then the strict transform of $Y$ inside $\Bl_Z(X)$ lies inside the affine blowup, namely
\begin{equation*}
    Y^{\st} \hookrightarrow \Bl^{\circ}_{Z}(X).
\end{equation*}
Indeed, it is enough to show that $Y^{\st}$ does not intersect $X^{\st}_{\zeta}$. Using the properties from \S \ref{section: blowups}, \ref{section: strict transforms}, this follows from
\begin{equation*}
  Y^{\st} \cap  X^{\st}_{\zeta} = \Bl_{Z}(Y) \cap \Bl_Z(X_{\zeta}) = \Bl_{Z}(Y \cap X_{\zeta}) = \Bl_{Z}(Z) = \emptyset.
\end{equation*}


\subsubsection{Deformations centered at a set of points.}\label{subsubsection: deformations centered at a set of points}
We will often perform the affine blowup construction over a set of closed points of $\Gm \hookrightarrow \A^1$. Namely, suppose we have a scheme $X$ and consider $X \times \Gm$. Suppose we are given a set $(z_i \mid i \in I)$ of pairwise distinct closed points of $\Gm$ and a set of closed subschemes $\eZ = (Z_i \subseteq X \mid i \in I)$. We regard $Z_i = Z_i \times \{z_ i\} \hookrightarrow X \times \Gm$ as a closed subscheme.

We the may blow up independently each $Z_i$ inside $X \times \Gm$; we denote the result by
\begin{equation*}
    N_{q=z_i}(Z_{z_i}, X) := \Bl^{\circ}_{Z_{i} \times \{z_i\}}(X \times \Gm^{\rot})
\end{equation*}
where the subscript $q=z_i$ emphasizes that the deformation is centered at the point $z_i$ of $\Gm = \Spec(k[q^{\pm 1}])$. 

We denote the associated gluing of deformations over the set of points $(z_i \mid i \in I)$ by
\begin{equation*}
    \eN_{I}(\eZ, X) := \Bl^{\circ}_{\eZ}(X).
\end{equation*}
This is defined as the affine blowup of $X \times \A^1$ centered in $\coprod_{i \in I} Z_{i} \times \{z_i\} \hookrightarrow X \times \A^1$. However, note the when $I$ is infinite, the center of the blowup is merely an ind-scheme. 

Explicitly, if $I$ is finite, this gluing is given by inductively blowing up at $Z_i \times \{ z_i \}$ in some order; the order does not matter as these are disjoint. For infinite $I$, one takes the limit over all finite subsets of $I$ -- these are schemes, hence their limit also is.


\subsection{Reductive groups}
\subsubsection{Reductive groups.}\label{subsubsection: reductive groups}
Let $G$ be a connected reductive group over an algebraically closed field $k$. 
Denote by $T$ its universal maximal torus; we also write $T = T_G$ whenever we need to stress the ambient group. For the sake of concreteness, we fix an embedding $T \leq G$. We write $X_{\bullet} = X_{\bullet}(T)$ resp. $X^{\bullet}=X^{\bullet}(T)$ for the cocharacter resp. character lattice of $T$. There is a natural identification $T(k) = X_{\bullet}(T) \otimes _{\Z} k^{\times}$.

We choose a Borel $B \leq G$ containing $T$ and denote by $X^{+}_{\bullet} = X^{+}_{\bullet}(G)$ the set of dominant cocharacters with respect to $B$.
%
We denote $B' \leq G$ the Borel subgroup opposite to $B$; all parabolics and parahorics will be taken with respect to $B'$.

\subsubsection{Chevalley forms.}
We use the same notation as in \S \ref{subsubsection: reductive groups} over the ring of coefficients $\k$ for our cohomology theories; for example we may take $\k = \Z$. In particular, we denote by $G := G_{\k}$ the Chevalley form over $\k$. We further denote by $T := T_{\k}$ its universal maximal torus, $B := B_{\k}$ a chosen Borel subgroup and so on. We omit the subscript $\k$ when there is no risk of confusion. All reductive groups in this paper are split.


\subsubsection{Root data.}
We write $\Phi_{\bullet} \subseteq X_{\bullet}$ resp. $\Phi^{\bullet} \subseteq X^{\bullet}$ for the coroots resp. roots of $G$. They respectively generate the coroot lattice $Q_{\bullet} := \Z\Phi_{\bullet} \leq X_{\bullet}$ and root lattice $Q^{\bullet} :=\Z\Phi^{\bullet} \leq X^{\bullet}$.
Altogether, $(X_{\bullet}, \Phi_{\bullet}, X^{\bullet}, \Phi^{\bullet})$ denotes the root datum of $G$; its automorphism group is the Weyl group $W$.
%
Whenever the ambient reductive group needs to be emphasized, we also write $X_{\bullet, G}$, $\Phi_{\bullet, G} $, $X^{\bullet}_G$, $\Phi^{\bullet}_G$, $W_G$ and so on.

%
We write $\langle -, - \rangle$ for the natural pairing $X_{\bullet} \times X^{\bullet} \to \Z$. We denote by $(-)^\vee: \Phi^{\bullet} \to \Phi_{\bullet}$, $\beta \mapsto \widecheck{\beta}$ the unique bijection such that $\langle \widecheck{\beta}, \beta \rangle = 2$ for each root $\beta \in \Phi^{\bullet}$. Denote further $\Phi_{\bullet}^{+}$ resp. $\Phi^{\bullet}_{+}$ the sets of positive coroots and positive roots with respect to $B$. We let $\Phi_{\bullet}^{\sr} = \{\widecheck{\beta}_1, \dots, \widecheck{\beta}_d \}$ resp. $\Phi^{\bullet}_{\sr} = \{\beta_1, \dots, \beta_d \}$ denote the subsets of simple coroots and simple roots. 

We write $\leq$ for the usual partial order on $X_{\bullet}$. By definition, $\lambda' \leq \lambda$ if and only if $\lambda - \lambda' \in \Z_{\geq 0} \Phi^{+}_{\bullet}$. Similarly in the dual case of $X^{\bullet}$.

We denote $X_{\bullet,\Q} := X^{\bullet}(T) \otimes_{\Z} \Q$ and $X^{\bullet}_{\Q} := X^{\bullet}(T) \otimes_{\Z} \Q$.



\subsubsection{Langlands dual and its representations.}\label{section: lnglands dual and its representations}
If $G$ is a reductive group with a root datum $(X_{\bullet}, \Phi_{\bullet}, X^{\bullet}, \Phi^{\bullet})$, its Langlands dual is the reductive group $\widecheck{G}$ with the flipped root datum given by $(\widecheck{X}_{\bullet}, \widecheck{\Phi}_{\bullet}, \widecheck{X}^{\bullet}, \widecheck{\Phi}^{\bullet}) := (X^{\bullet}, \Phi^{\bullet}, X_{\bullet}, \Phi_{\bullet})$.

Given any $\mu \in X^{+}_{\bullet} = \widecheck{X}^{\bullet}_+$, we denote $V_{\mu}$ the corresponding irreducible representation of $\widecheck{G}$ of highest weight $\mu$. 
We further denote 
\begin{equation}\label{notation: weights in mu}
    \wt(\mu) := W \cdot \{ \lambda \in X^{+}_{\bullet} \mid \lambda \leq \mu \} \subseteq X_{\bullet}
\end{equation}
the union of $W$-orbits of elements $\lambda \in X^{+}_{\bullet}$ with $\lambda \leq \mu$. In terms of $\widecheck{G}$, it is the subset of characters of nonzero weight spaces in $V_{\mu}$.


\subsubsection{Weyl groups.}
Denote $W = N_G(T)/T$ the Weyl group of $T$ in $G$. Equivalently, it is the Weyl group of the root system $(X_{\bullet}, \Phi_{\bullet}, X^{\bullet}, \Phi^{\bullet})$. It acts on $T$ and all the above objects in the standard way. 

Given any root $\beta \in \Phi^{\bullet}$, we denote the corresponding reflection
\begin{equation*}
    s_{\beta}: X_{\bullet} \to X_{\bullet}, \qquad \lambda \mapsto \lambda - \langle \beta, \lambda \rangle \widecheck{\beta}.
\end{equation*}
It induces an automorphism  of $T$, in multiplicative notation given by
\begin{equation*}
s_{\beta}: T \to T, \qquad s_{\beta}(c) = c \cdot \widecheck{\beta}(\beta(c))^{-1}.    
\end{equation*} 
The reflections $s_{\beta}$ induce elements of the Weyl group $W$. Moreover, $W$ has the structure of a Coxeter group, generated by the simple reflections $s_{\beta}$ for $\beta \in \Phi^{\bullet}_{s}$. 



\subsubsection{Simply connected covers and adjoint groups.}\label{subsubsection: simply connected covers and adjoint groups}
Denote $G^{\der}$ the derived group of $G$. This is a semisimple group. It is often useful to consider the finitely many (semisimple) groups isogeneous to $G^{\der}$. 
In particular, let $G^{\sc}$ be the simply connected cover of $G^{\der}$; its maximal torus is denoted $T^{\sc}$. Similarly, let $G^{\ad}$ be the adjoint group of $G^{\der}$, with maximal torus $T^{\ad}$. The adjoin group is given directly by the quotient of $G$ by its center: $G^{\ad} = G/Z(G)$. Concisely, we have the diagram
\begin{equation*}
\begin{tikzcd}
G^{\sc} \arrow[d] \arrow[rd]\\
G^{\der} \arrow[r, hook] \arrow[d] & G \arrow[ld, two heads]\\
G^{\ad}
\end{tikzcd}    
\end{equation*}
The groups in the left column are all semisimple; the vertical maps between them are central isogenies. Each connected reductive group isogenous to $G^{\der}$ sits in between $G^{\sc}$ and $G^{\ad}$.
When $G=G^{\sc}$, we say $G$ is simply connected group (implicitly assuming it is reductive); similarly when $G=G^{\ad}$ is adjoint.

The above diagram also induces corresponding maps on the universal maximal tori, and consequently on other root theoretic data. In particular, these maps naturally identify the root resp. coroot lattices $\Z\Phi^{\bullet}$ resp. $\Z\Phi_{\bullet}$ of all the groups; we will implicitly use these identifications. 

We write $X^{\sc}_{\bullet} = X_{\bullet}(T^{\sc})$ and $X_{\sc}^{\bullet} = X^{\bullet}(T^{\sc})$ for the cocharacter and character lattice of $T^{\sc}$, as well as $X_{\bullet}^{\ad} = X_{\bullet}(T^{\ad})$ and $X_{\ad}^{\bullet} = X^{\bullet}(T^{\ad})$ for the cocharacter and character lattices of $T^{\ad}$, and so on.
For any cocharacter $\lambda \in X_{\bullet}$ of $T$, we write $\lambda^{\ad} \in X^{\ad}_{\bullet}$ for the cocharacter of $T^{\ad}$ given by the image of $\lambda$ under the natural projection $T \twoheadrightarrow T^{\ad}$. 



\subsubsection{Algebraic fundamental group.}\label{section: borovoi fundamental group}
Denote
\begin{equation}
 \pi_1(G) := X_{\bullet} / X^{\sc}_{\bullet} = X_{\bullet} / \Z\Phi_{\bullet}  
\end{equation}
the \textit{algebraic fundamental group of $G$}, see \cite[Definition 1.3]{Bor98}. It recovers the usual fundamental group of the topological space $G(\C)$ over $\C$. We have $G=G^{\sc}$ if and only if $\pi_1(G)=1$. The natural action of $W$ on $\pi_1(G)$ is trivial and the fundamental group $\pi_1(G)$ is functorial in $G$; see \cite[Lemma 1.2 and its proof]{Bor98}.
%
Given $\lambda \in X_{\bullet}$, we denote by $\overline{\lambda} \in \pi_1(G)$ the corresponding class.


We choose a family of representatives for elements of the fundamental group under the surjection $X_{\bullet}(T) \twoheadrightarrow \pi_1(G)$: given an element $\sigma \in \pi_1(G)$, let $ \chi(\sigma) \in X_{\bullet}(T)$ be the chosen representative. 


\subsubsection{Shifted actions.}\label{section: shifted actions}
Given $\alpha \in X_{\bullet}$, we denote $\bullet_{\alpha}$ the $\alpha$-shifted action
\begin{align*}
    (W, \bullet_{\alpha}) & \acts X_{\bullet} \\
    w \bullet_{\alpha} \mu & = w(\mu - \alpha) + \alpha, \qquad \forall \ w \in W.
\end{align*}
It identifies with the natural action $W \acts \alpha + X_{\bullet}$ of $W$ on the $\alpha$-shifted copy of $X_{\bullet}$.
When $\alpha = \rho$, this is the \emph{dot action}. We use the same notation for the induced action on $T$.

This notation naturally extends to the following setup. Assume that $\omega \in X^{\ad}_{\bullet}$ is a cocharacter of the adjoint group $G^{\ad}$. We then pick any lift $\alpha \in X_{\bullet, \Q}$ of $\omega$ under the natural projection $X_{\bullet, \Q} \twoheadrightarrow X^{\ad}_{\bullet, \Q}$. 
The associated $\alpha$-shifted action $(W, \bullet_{\alpha}) \acts X_{\bullet, \Q}$ on the vector space $X_{\bullet, \Q}$, given by $w \bullet_{\alpha} \mu = w(\mu - \alpha) + \alpha = w(\mu) - w(\alpha) + \alpha$, preserves the subset $X_{\bullet}$ and is independent of the choice of the lift $\alpha$. 
%
Indeed, it is enough to check that for any $w \in W$, the element $\alpha - w(\alpha)$ lies in $X_{\bullet}$ and is independent of the choice of $\alpha$. By induction on the length of $w$, both statements may be checked on simple reflections, where they are clear since $\alpha - s_{\beta}(\alpha) = \langle \beta, \alpha \rangle \widecheck{\beta} = \langle \beta, \omega \rangle \widecheck{\beta}$.
%

Altogether, given $\omega \in X_{\bullet}^{\ad}$, we obtain a well-defined action
\begin{equation*}\label{equation: action shifted by adjoint cocharacter}
    (W, \bullet_{\omega}) \acts X_{\bullet}. 
\end{equation*}
given by the same formula with respect to any lift of $\omega$.

%; it identifies with the natural action $W \acts \lambda + X_{\bullet}$ on the shifted copy of $X_{\bullet}$


\subsubsection{Levi subgroups.}\label{section: levi subgroups}
Given a cocharacter $\lambda: \Gm \to G$, we denote
$L_{\lambda}$ the centralizer of the image of $\lambda$ inside $G$. 
Then $L_{\lambda} \leq G$ is a Levi subgroup; it is in particular connected reductive. We denote by $P_{\lambda} \leq G$ the corresponding parabolic subgroup with respect to $B'$ (containing $L_{\lambda}$ and $B'$). See \cite[\S 21.i and Theorem 13.33]{Mil17}.

We call $L_{\lambda} \leq P_{\lambda}$ the Levi and parabolic associated to $\lambda$; note that this construction depends on the ambient reductive group $G$. %Later in our applications, it will be some Levi subgroup to start with.
%
Note that $T \leq L_{\lambda}$ is a maximal torus and $X_{\bullet} = X_{\bullet}(T)$ is the cocharacter lattice of $L$. The intersections of $B$ resp. $B'$ with $L$ define a pair of opposite Borel subgroups $B_L$ resp. $B'_L$ of $L$. 

For any reductive group $G$, the Levi subgroups and parabolic subgroups of $G$ and $G^{\ad}$ are in natural bijection. Consequently, we may consider the subgroups $L_{\omega}$, $P_{\omega}$ of $G$ labeled by $\omega \in X^{\ad}_{\bullet}$.
%In other words, the set of isomorphism classes of Levi resp. parabolic subgroups only depends on the central isogeny class of $G$


\subsubsection{Coordinates and invariant polynomials.}\label{subsubsection: coordinates and invariant polynomials}
Let $G$ be an affine group scheme over a field $k$.
The representation ring of $G$ is denoted $R(G)$. Given a coefficient ring $\k$, its $\k$-linearization is denoted $R(G;\k):= R(G) \otimes_{\Z} \k$.




Assume $G = G$ is a split reductive group over $k$ of rank $n$ with maximal torus $T$. 
We then have its Chevalley form $G_{\k}$ over $\k$, and we denote $T_{\k}$ is maximal torus. The representation ring of $G$ is given by
\begin{equation*}
    R(G) = \O(T_{\Z} \sslash W) \qquad \text{and} \qquad R(G; \k) = \O(T_{\k} \sslash W).
\end{equation*}
As a module, $R(G)$ is generated by classes of irreducible representations of $G$, denoted $V_{\mu}$ for $\mu \in X^{\bullet}_{+}$. For ring generators, one may take the $[V_{\mu_j}]$ with $\mu_j$ for $j =1, \dots, n$, where $V_j$ are chosen generators of $X^{\bullet}_{+}$. We denote these $n$ elements by
\begin{equation*}
    e_j := e_j^{G} := [V_{\mu_j}] \qquad j = 1, \dots, n.
\end{equation*}
In particular:
\begin{itemize}
\item 
If $G=T$ to start with, we denote these coordinates by
$\O(T_{\Z}) = \Z[X^{\bullet}] = \Z[t_1^{\pm 1}, \dots, t_n^{\pm 1}]$.
In a coordinate free way, we also denote by $t_{\nu}$ the coordinate coming from the character $\nu$. 
%
\item
If $G = GL_n$, we have $\O(T_{\Z} \sslash W) = \Z[e_1, \dots, e_{n-1}, e_n^{\pm 1}]$. Here, $e_j$ is the $j$-th elementary symmetric polynomial on $T_{\Z}$.
%
\item
For general $G$, interpreting any $V_{\mu}$ as a group morphism $\rho_{\mu}: G \to \GL_n$ for a suitable $n$, we get the above classes as pullbacks of the standard representation: $e^{G}_j = r^*_{\mu_j} (e_1)$.
\item 
For $G$ simply connected, we take $e_j := [V_{\omega_j}]$ the classes of the fundamental representations with highest weights $\omega_j$, for $j = 1, \dots, n$.
\end{itemize}


\begin{remark}\label{remark: representation ring in the simply conneted case}
For reductive $G$, the ring $\O(T_{\Z} \sslash W)$ may be subtle. Namely, when $\pi_1(G) \neq 1$, it will usually acquire finite quotient singularities -- see \cite[\S 7]{BZ00} for the case of $G = \PGL_3$.

On the other hand, for $G$ simply connected, the above ring is a polynomial ring: $\O(T_{\Z} \sslash W) = \Z[e_1, \dots, e_n]$ with distinguished coordinates given by the fundamental representations $e_j \in \O(T \sslash W)$.
\end{remark}


\subsubsection{Lie algebras.}
Given a split reductive group $G$ over $\k$, we denote its Lie algebra by $\g$. In particular, the Lie algebra of its maximal torus $T$ is denoted $\t$. When $G$ is simply connected, we have $T \sslash W = \t \sslash W$, an affine space from Remark \ref{remark: representation ring in the simply conneted case}.




\subsection{Loop groups and affine flag varieties}\label{section: loop groups and affine flag varieties}

\subsubsection{Affine root data.}
We denote by $\Phi^{\bullet}_{\aff} = \Phi^{\bullet} \oplus \Z\hbar$ the set of {\it affine real roots} associated to the root system of $G$. We write elements of this set as $\beta + m \hbar$ with $\beta \in \Phi^{\bullet}$ and $m \in \Z$.

\subsubsection{Affine and extended Weyl groups.}
Let $\Z\Phi_{\bullet} \leq X_{\bullet}$ be the coroot lattice inside the cocharacter lattice.
We denote by
\begin{equation*}
W_{\aff} := \Z\Phi_{\bullet} \rtimes W
\qquad
\text{and}
\qquad
W_{\ext} := X_{\bullet} \rtimes W
\end{equation*}
the \textit{affine Weyl group} and \textit{extended affine Weyl group} associated to the root system of $G$.
These semi-direct products are constructed from the natural action of $W$ on $\Phi_{\bullet}$ resp. $X_{\bullet}$. They thus have preferred splittings, allowing us to regard $W$ as a subgroup.
Clearly, $W_{\aff} \leq W_{\ext}$ is a normal subgroup; the quotient identifies as $W_{\ext} / W_{\aff} = \pi_1(G)$. In particular $W_{\ext} = W_{\aff}$ if $G$ is simply connected.
%
Also note that $W_{\aff}$ only depends on the isogeny class of $G$, since this is the case for both $W$ and $\Phi_{\bullet}$. On the other hand, $W_{\ext}$ depends on the actual group $G$.


The image of a coroot $\widecheck{\beta} \in \Phi_{\bullet}$ is denoted $\tau_{\widecheck{\beta}}:= (\widecheck{\beta}, 1) \in W_{\aff}$. We also have the finite reflections $s_{\beta} \in W \leq W_{\aff}$ labeled by $\beta \in \Phi^{\bullet}$.
%
The affine Weyl group $W_{\aff}$ has the structure of a Coxeter group, with simple reflections given by
\begin{equation*}
    s_{\beta, m} := (\tau_{m\widecheck{\beta}}, s_{\beta}) \in \Z\Phi_{\bullet} \rtimes W, \qquad \beta \in \Phi_{\bullet}, \ m \in \Z. 
\end{equation*}
Similar discussion applies to $W_{\ext}$.


\subsubsection{Shifted affine Weyl group actions.}\label{section: shifted affine weyl group actions}
Consider the following family of actions of $W_{\aff}$ on $X_{\bullet}$.
Let $\ell \in \N$ be a natural number and $\alpha \in X_{\bullet}$ any cocharacter. Then we define the \emph{$\ell$-dilated $\alpha$-shifted action}
\begin{align*}
    (W_{\aff}, \bullet^{\ell}_{\alpha}) & \acts X_{\bullet} \\
    (\lambda, w ) \bullet^{\ell}_{\alpha} \mu & = w(\mu - \alpha - \ell \lambda) + \alpha, \qquad \forall \lambda \in \Z\Phi_{\bullet}, \ w \in W
\end{align*}

For $\ell = 1$ and $\alpha$ trivial, this is the natural of $W_{\aff}$ on $X_{\bullet}$. For general $\ell$, it is the \emph{$\ell$-dilated action}. For $\alpha = \rho$ and general $\ell$, we get the \emph{$\ell$-dilated dot action}. We use the same notation for the induced action on $T$.


Given any cocharacter $\omega \in X^{\ad}_{\bullet}$ of the adjoint group $G^{\ad}$, we obtain the $\omega$-shifted action
\begin{equation*}
   (W_{\aff}, \bullet^{\ell}_{\omega})  \acts X_{\bullet} 
\end{equation*}
given by the same formula with respect to any representative of $\omega$ inside $X_{\bullet, \Q}$. It is well-defined by the argument from \S \ref{section: shifted actions}.







\subsubsection{Loop rotation torus.}
We denote by $\Gm^{\rot}$ the one dimensional torus with coordinate ring $\O(\Gm^{\rot}) = k[q^{\pm 1}]$ and call it the \textit{loop rotation torus}. It plays a special role, so we include the superscript $\rot$ to distinguish from other tori. We denote its closed points by $\zeta \in \Gm^{\rot}$.

For $\ell \in \N$, we write $\mu_{\ell} \leq \Gm^{\rot}$ for the subgroup scheme of order $\ell$ roots of unity. We write $\mu_{\infty} := \colim_{\ell \in \N} \mu_{\ell} \leq \Gm^{\rot}$ for the ind-subgroup scheme of all roots of unity. This gives a decomposition $\Gm^{\rot} = (\Gm^{\rot} - \mu_{\infty}) \sqcup \mu_{\infty}$.


A closed point $\zeta \in \Gm^{\rot}$ is called {\it generic} if it is not a root of unity, i.e. if it is torsionfree element of the underlying multiplicative group. Otherwise, $\zeta \in \mu_{\infty}$ is a root of unity.
For any $\zeta \in \Gm^{\rot}$, we write
\begin{equation*}
\ord(\zeta) =
\begin{cases}
\ell, \qquad \text{if} \ \zeta \in \mu_{\infty} \ \text{is a root of unity of primitive order} \ \ell \in \N, \\
0, \qquad \text{if} \ \zeta \in \Gm^{\rot} - \mu_{\infty} \ \text{is generic}.
\end{cases}
\end{equation*}
We write $\Sigma_{\ell} \hookrightarrow \mu_{\ell}$ for the cyclotomic subscheme of roots of unity of primitive order $\ell$. 

Denoting $\hbar$ the real affine root in the Lie algebra sense, our group coordinate $q$ is given as $q = e^{\hbar}$. Since we work multiplicatively throughout, we do not need the exponential notation.


\subsubsection{Loop groups with loop rotation.}
Given a connected reductive group $G$ over $k$, we denote by $\lo G$ its loop group and by $\lop G$ its positive loop group, with functors of points
\begin{equation*}
    \lo G: R \mapsto G(R\llp t \rrp) \qquad \text{and} \qquad \lop G: R \mapsto G(R\llb t \rrb).
\end{equation*}
We have $\lop G \leq \lo G$. The group of connected components of the loop group $\lo G$ identifies with $\pi_1(G)$, while the positive loop group $\lop G$ is connected. We have an embedding $G \leq \lop G \leq \lo G$ of the original reductive group $G$ via constant loops.


The torus $\Gm^{\rot}$ acts on $\lo G$ by the loop rotation 
\begin{align*}
\Gm^{\rot} \times \lo G \to \lo G  , \\
(\zeta, g(t)) \mapsto g(\zeta t).
\end{align*}
This action clearly restricts to $\lop G$. 
%
The extended loop groups are given by the semidirect products along loop rotation
\begin{equation*}
\lop G \rtimes \Gm^{\rot} \leq \lo G \rtimes \Gm^{\rot}.   
\end{equation*}
These semidirect products become trivial after restriction to the constant loops $G \leq \lop G$.

Given $\ell \in \N$, we write $\lo_{\ell} G \leq \lo G$ and $\lopl G \leq \lop G$ the $\ell$-dilated subgroups; they represent the subfunctors
\begin{equation*}
    \lol G: R \mapsto G(R\llp t^{\ell} \rrp) \qquad \text{and} \qquad \lopl G: R \mapsto G(R\llb t^{\ell} \rrb).
\end{equation*}
While abstractly isomorphic to the usual loop groups, the geometry of their embeddings is useful.
%
For $\ell = 0$, we denote by convention $\lo_0 G = \mathrm{L^+_{0}} G = G$, embedded in $\lo G$ as constant loops.


\subsubsection{Extended tori.}\label{subsubsection: extended tori}
In particular, we have the extended torus $T \times \Gm^{\rot} \leq \lo G \rtimes \Gm^{\rot}$ embedded via constant loops.
Many computations will take place within this ambient extended loop group.
For notational convenience, we pick coordinates
\begin{equation*}
    \O(T\times \Gm^{\rot}) = k[t_1^{\pm 1}, \dots, t_n^{\pm 1}][q^{\pm 1}].
\end{equation*}


Any character of $X^{\bullet}(T)$, and in particular any root $ \beta \in \Phi^{\bullet} \leq X^{\bullet}(T)$, can be tautologically regarded as a function $\beta \in \O(T)$ given by the composition $\iota \circ \beta: T \to \Gm \to \A^1$, where $\iota: \Gm^{\rot} \hookrightarrow \A^1$ is the obvious inclusion.
%
Similarly, any character of $X^{\bullet}(T\times \Gm^{\rot})$, and in particular any affine root $\beta + m \hbar$, can be tautologically regarded as a function $\beta + m\hbar \in \O(T \times \Gm^{\rot})$ given by the composition $\iota \circ (\beta + m\hbar): T \times \Gm^{\rot} \to \Gm \to \A^1$.
%
It is sometimes customary to denote these functions as
\begin{equation*}
   e^{\beta} = \beta \in \O(T) \qquad \text{and} \qquad e^{\beta + m\hbar} = e^{\beta} \cdot e^{m\hbar} = {\beta} \cdot q^m \in \O(T \times \Gm^{\rot}).
\end{equation*}
Since we are working multiplicatively, we will avoid this exponential notation.



On the other hand, given any cocharacter $\lambda \in X_{\bullet}(T)$, regarded as a homomorphism $\lambda: \Gm \to T$, we denote $t^{\lambda} := \lambda(t) \in \lo T$ the image of the formal variable $t$ under the map $\lo \lambda: \lo \Gm \to \lo T$. More generally, given any element $a \in k \llb t \rrb$, we denote $a^{\lambda} := \lambda(a) \in \lo T$ its image under $\lo \lambda$.



\subsubsection{Root spaces in loop groups.}\label{paragraph: root spaces in loop groups}
Given an affine root $\beta + m\hbar \in \Phi^{\bullet} \oplus \Z \hbar$, we get the homomorphism $\Ga \to \lo G$ given by $x \mapsto u_\beta(t^mx)$ and denote $U_{\beta + m\hbar} \leq \lo G$ the corresponding root subgroup of $\lo G$ given by its image; see \cite[\S 4.3]{RW22}.

\begin{lemma}\label{lemma: weights of extended torus on affine root groups}
The adjoint action of $T \times \Gm^{\rot}$ restricts to each $U_{\beta + m\hbar}$, where it has weight
\begin{equation*}
(x, \zeta) \mapsto \beta(x) \cdot \zeta^m.  
\end{equation*}
\end{lemma}
\begin{proof}
Given any point $(x, \zeta) \in T \times \Gm^{\rot}$, compute
\begin{align*}
   (x, \zeta) \cdot (u_{\beta}(z t^m), 1) \cdot (x, \zeta)^{-1} 
   = (x, 1) \cdot (u_{\beta}(\zeta^m \cdot zt^m), 1) \cdot (x^{-1}, 1)
   = (u_{\beta}(\beta(x) \zeta^m \cdot z t^m), 1).
\end{align*}
\end{proof}

\begin{lemma}\label{lemma: conjugation of affine root spaces}
For any $\lambda \in X_{\bullet}$ we have 
\begin{equation*}
 t^{-\lambda} \cdot U_{\beta + m\hbar} \cdot t^{\lambda} = U_{\beta + (m- \langle \lambda, \beta \rangle)\hbar}
\end{equation*}
\end{lemma}
\begin{proof}
    See \cite[Lemma 4.3. (2)]{RW22}.
\end{proof}





\subsubsection{Parahoric subgroups.}\label{subsubsection: parahoric subgroups}
The preimage of the Borel subgroup $B' \leq G$ under the evaluation map $\lop G \to G$ gives rise to the standard Iwahori subgroup $\Iw \leq \lop G$. Given any proper subset $J$ of the extended Dynkin diagram of $G$, we denote the associated standard parahoric subgroup $\eP_J$. We have $\Iw \leq \eP_J \leq \lop G$. To avoid too many subscripts, we sometimes put $J$ in the superscript and write $\eP^J := \eP_J$.

A parahoric subgroup $\eP_{J}$ sits in a canonical short exact sequence
\begin{equation*}
1 \to \eP_{J}^{\uni} \to \eP_{J} \to L_{J} \to 1    
\end{equation*}
where $\eP_{J}$ is its \textit{pro-unipotent radical} and $L_{J}$ is its reductive \textit{Levi quotient}. 
For all questions regarding equivariant topological invariants, the difference between $\eP_{J}$ and $L_{J}$ is irrelevant.

For more details on parahoric subgroups, see \cite[\S 2.1.5]{Yun15} and references there.

\subsubsection{Affine flag varieties.}
Given a parahoric subgroup $\eP \leq \lop G$, the associated affine flag variety is the fpqc quotient
\begin{equation*}
    \Fl_{\eP} = \lo G / \eP = \lo G \rtimes \Gm^{\rot} / \eP \rtimes \Gm^{\rot}.
\end{equation*}
It is an ind-projective ind-scheme. It carries a left transitive action of $\lo G \rtimes \Gm^{\rot}$; we often restrict the above action to various subgroups of $\lop G \rtimes \Gm^{\rot}$, such as $T \times \Gm^{\rot}$.

When $\eP= \eP_J$ is the standard parahoric, we also write $\Fl^J := \Fl_{\eP_J}$.
For $\eP = \lop G$, we have the affine Grassmannian
\begin{equation*}
    \Gr_{G} = \lo G / \lop G = \lo G \rtimes \Gm^{\rot} / \lop G \rtimes \Gm^{\rot}.
\end{equation*}

For each $\lambda \in X_{\bullet}$, the element $t^{\lambda}$ gives a distinguished point on $\Fl_{\eP}$ and in particular on $\Gr_G$. These are precisely the fixed-points under the whole torus $T \times \Gm^{\rot}$. We call them the \emph{coordinate fixed-points}.


\subsubsection{Local Hecke stacks.}\label{section: local quantum hecke stacks}
Taking the quotient by the above action, we obtain
\begin{equation*}
   \lop G \rtimes \Gm^{\rot} \backslash \Fl_{\eP} = \lop G \rtimes \Gm^{\rot} \backslash \lo G / \eP = \lop G \rtimes \Gm^{\rot} \backslash \lo G \rtimes \Gm^{\rot} / \eP \rtimes \Gm^{\rot}.
\end{equation*}
In the case of affine Grassmannian, we call this quotient the \emph{quantum Hecke stack}
\begin{equation*}
   \lop G \rtimes \Gm^{\rot} \backslash \Gr_{G} = \lop G \rtimes \Gm^{\rot} \backslash \lo G / \lop G = \lop G \rtimes \Gm^{\rot} \backslash \lo G \rtimes \Gm^{\rot} / \lop G \rtimes \Gm^{\rot}.
\end{equation*}



\subsubsection{Affine Schubert varieties.}
The $\lop G \rtimes \Gm^{\rot}$ orbits on $\Gr_G$ are the affine Schubert cells $\Gr_{\mu}$, equipped with the reduced scheme structure. By the Iwasawa decomposition, they are labeled by the discrete set $X^{+}_{\bullet}$. The affine Schubert cell $\Gr_{\mu}$ is an affine bundle over the corresponding partial flag variety $G/P_{\mu}$.

The affine Schubert varieties $\Gr_{\leq \mu}$ are the closures of the affine Schubert cells, equipped with the reduced scheme structure. They have a locally closed decomposition into a disjoint union of the affine Schubert cells
\begin{equation*}
    \Gr_{\leq \mu} = \bigsqcup_{\mu' \leq \mu} \Gr_{\mu'}.
\end{equation*}
Each $\Gr_{\leq \mu}$ is a singular projective variety, acted on by $\lop G \rtimes \Gm^{\rot}$. We refer to \cite{Zhu15} for an excellent introduction.

\begin{remark}
The affine Grassmannian $\Gr_{G}$ has a natural affine paving. Indeed, we have an action $\Iw \acts \Gr_{G}$ of the Iwahori group scheme. Its fixed points are precisely $t^{\lambda}$ with $\lambda \in X_{\bullet}$. The action of $\Iw$ on $t^{\lambda}$ factors through a finite-dimensional quotient. Since the reductive quotient $T$ of $\Iw$ stabilizes $t^{\lambda}$, its orbit under $\Iw$ is an affine space. Consequently, we deduce that Iwahori orbits provide an affine paving of $\Gr_G$, and in particular of each affine Schubert variety $\Gr_{\leq \mu}$. 

Also see \cite[(5.2)]{HHH05} and references there.
\end{remark}




\subsubsection{Dilated variants.}\label{subsubsection: dilated variants}
We will use the above constructions in the \textit{$\ell$-dilated} context. Also, the ambient reductive group will in fact be a suitable pseudo-Levi subgroup $L$ of $G$. We spell this out in detail.

Let $\lambda \in X_{\bullet} = X_{\bullet}(T)$. Consider the subset $I_{\ell, \lambda}$ of affine reflections fixing $\lambda$. Denote $\eP^{\lambda}_{L, \ell} \leq \lo_{\ell} L$ the standard parahoric subgroup corresponding to $I_{\ell, \lambda}$ with respect to the opposite Borel $B'$ -- up to passage to the opposite Borel, see \cite[\S 2.1.4]{BBSV22}).
In terms of Bruhat--Tits theory, $\eP^{\lambda}_{L, \ell}$ is the parahoric subgroup of $\lo_{\ell} L$ corresponding to the facet $\mathrm{f}$ containing $\lambda$, see \cite{RW22}. 
Explicitly, by \cite[equation (2.8)]{BBSV22} together with our opposite Borel convention, it is given as
\begin{equation*}
    \eP^{\lambda}_{L, \ell} = \lo_{\ell} L \cap t^{\lambda} \cdot \lop L \cdot t^{-\lambda}.
\end{equation*}
See \cite[\S 4.2 -- \S 4.3]{RW22} and \cite[\S 2.1]{BBSV22}.
%
We denote 
\begin{equation*}
\Fl^{\lambda}_{L, \ell} := \lo_{\ell} L / \eP^{\lambda}_{L, \ell}    
\end{equation*}
the corresponding partial affine flag variety of the $\ell$-dilated loop group of $L$.
Note that these constructions depend only on the stabilizer of $\lambda$ in $W_{\ell, \aff}$.


\subsubsection{Functoriality.}
The constructions in \S \ref{section: loop groups and affine flag varieties} are functorial in group homomorphisms $G \to G'$ of reductive groups. Specifically, we have an induced map $\Gr_{G} \to \Gr_{G'}$ compatible with the equivariance, Schubert stratifications, and so on. Consequently, many constructions can be reduced to the case of $GL_n$.



\subsection{Root subsystems and centralizers}\label{section: root subsystems and centralizers}
In order to effectively manipulate centralizers of semisimple elements inside loop groups, we recall some technical material on subsystems of finite and affine root systems. We then briefly recall the structure of such centralizers. In the finite case, these are called pseudo-Levi subgroups. 
We work over an algebraically closed field $k$. For simplicity, we assume $k$ is of characteristic zero.

\subsubsection{Closed subsystems of root systems.}

By an \emph{abstract root system} $\Psi^{\bullet}$ we mean a root system of a Kac--Moody algebra $\g_{\Psi^{\bullet}}$ in the sense of \cite{Kac83}; the notations below make sense in this generality.

\begin{setup}\label{setup: finite and affine root systems}
For our purposes, we will assume that $\Psi^{\bullet}_1$ is one of the following.
\begin{itemize}
    \item[(i)] $\Phi^{\bullet}_G$ is a finite root system, realized in a reductive group $G$ with respect to the chosen maximal torus $T \leq G$. We denote $\g$ its Lie algebra and $\t \subseteq \g$ the corresponding Cartan subalgebra.
    \item[(ii)] $\Phi^{\bullet}_{G, \aff} = \Phi^{\bullet}_G \oplus \Z\hbar$ is the real untwisted affine root system, realized by the affine real roots in the loop group $\lo G$ of $G$ with respect to the extended torus $T \times \Gm^{\rot}$. We denote $\g_{\aff}$ its Lie algebra and $\t_{\aff} = \t \oplus \Ga^{\rot}$ the corresponding Cartan subalgebra.
\end{itemize}
\end{setup}

Following \cite[Definition 2.4.1]{RV19}, we recall some terminology. Also see \cite[\S 1.2]{DL11a}.
\begin{notation}\label{notation: closed root subsystems}
Let $\Psi^{\bullet}_1$ be any ambient root system.
\begin{enumerate}
\item[(i)] 
A subset $\Psi^{\bullet}_2 \subseteq \Psi^{\bullet}_1$ is called a \emph{root subsystem} if it is itself a root system. 
Explicitly, this means $\Psi^{\bullet}_2$ is closed on reflections in its elements: $\forall \alpha, \beta \in \Psi^{\bullet}_2$, we have $s_{\alpha}(\beta) \in \Psi^{\bullet}_2$.
%
\item[(ii)] 
A subset $\Psi^{\bullet}_2 \subseteq \Psi^{\bullet}_1$ is called \emph{closed} if $\forall \alpha, \beta \in \Psi^{\bullet}_2$, also $ \alpha + \beta \in \Psi^{\bullet}_2$.
\end{enumerate}
A subset satisfying both (i) and (ii) is called a \emph{closed root subsystem}.
\end{notation}



\subsubsection{Projection and affinization.}\label{section: projection and affinization}
Affine and finite root systems compare through the projection map 
\begin{equation}\label{equation: projection of root systems}
\Phi^{\bullet}_{G, \aff} \twoheadrightarrow \Phi^{\bullet}_G, \qquad \beta + m\hbar \mapsto \beta.    
\end{equation}
Any subset $\Psi^{\bullet} \subseteq \Phi^{\bullet}_{G, \aff}$ can be projected to $\overline{\Psi}^{\bullet} \subseteq \Phi^{\bullet}_G$. If $\Psi^{\bullet}$ is a root subsystem, then also its projection is a root subsystem \cite[Definition 2.4.3]{RV19}. Moreover, if $\Psi^{\bullet} \subseteq \Phi^{\bullet}_{G, \aff}$ is a closed root subsystem, so is $\overline{\Psi}^{\bullet}$ by \cite[Proposition 4.1.1]{RV19}.

Vice versa, any subset $\Psi^{\bullet} \subseteq \Phi^{\bullet}_G$ can be lifted to a subset $ \widehat{\Psi}^{\bullet} \subseteq \Phi^{\bullet}_{G, \aff}$ by taking preimage \cite[Definition 2.4.4]{RV19}. This operation clearly preserves root subsystems, as well as the closed condition.



\subsubsection{Regular subalgebras.} 
From now on, we restrict attention to closed root subsystems; these parametrize a particularly useful class of subalgebras.
\begin{notation}
Let $\g_{\Psi^{\bullet}_1}$ be a Kac--Moody algebra. A subalgebra $\g' \subseteq \g_{\Psi^{\bullet}_1}$ is called \emph{regular} if it is stable under the adjoint action of some Cartan subalgebra $\t_{\Psi^{\bullet}_1}$ of $\g_{\Psi^{\bullet}_1}$. 
See \cite[Definition 12.1.3 and text below]{RV19} and \cite[\S 1]{FRT08}, as well as \cite[Chapter II, \S 2]{Dyn57} for the case of classical Lie algebras.
\end{notation}

Fix a Cartan subalgebra $\t_{\Psi^{\bullet}_1}$ of $\g_{\Psi^{\bullet}_1}$.
To a closed subroot system $\Psi^{\bullet}_2 \subseteq \Psi^{\bullet}_1$, we can associate the regular subalgebra $\g_{\Psi^{\bullet}_2}$ generated by the root spaces $\g_{\alpha}$, for $\alpha \in \Psi^{\bullet}_2$ as in \cite[Definition 12.1.3 and below]{RV19}. 
In the relevant cases, these give all regular subalgebras with respect to the chosen Cartan.


\begin{proposition}\label{proposition: closed root subsystems and regular subalgebras}
Assume that $\Psi^{\bullet}_1$ is either the finite root system $\Phi^{\bullet}_G$ or the affine real root system $\Phi^{\bullet}_{G, \aff}$ as in Setup \ref{setup: finite and affine root systems}. Denote $\t_{\Psi^{\bullet}_1} \subseteq \g_{\Psi^{\bullet}_1}$ the chosen Cartan subgroup in its Lie algebra. Then the above induces a bijection
    \begin{equation*}
    \{ \Psi^{\bullet}_2 \subseteq \Phi^{\bullet}_{\Psi^{\bullet}_1} \ \text{closed root subsystem}  \} \leftrightarrow \{ \g_{\Psi^{\bullet}_2} \subseteq \g_{\Psi^{\bullet}_1}  \text{regular subalgebra with respect to} \t_{\Psi^{\bullet}_1} \}.
\end{equation*}
\end{proposition}
\begin{proof}
The case of finite root systems is classical, going back to \cite{Dyn57}, \cite{BD49}. For the case of affine real root systems, see \cite[Corollary 12.1.5]{RV19}.
\end{proof}



Also see \cite[Theorem 12.3.1]{RV19} for other equivalent characterization of regular subalgebras of affine Kac--Moody algebras in terms of combinatorial data.

\begin{notation}
Since $k$ is algebraically closed of characteristic zero, to each regular subalgebra of $\g$ resp. $\g_{\aff}$ we can associate the corresponding connected subgroup of $G$ resp. $\lo G$.
We call these \emph{regular subgroups}.
\end{notation}

By construction, Proposition \ref{proposition: closed root subsystems and regular subalgebras} also parametrizes regular subgroups. The projection and affinization of closed root subystems from \S \ref{section: projection and affinization} induce corresponding constructions on regular subalgebras and regular subgroups.







\subsubsection{Centralizers of semisimple elements.}
Regular subalgebras and closed subroot systems in particular arise from centralizers of semisimple elements in the ambient Kac--Moody group. See \cite[Chapter II, \S 6, No.21, Corollaries to Theorem 6.7]{Dyn57} for the finite case.
\begin{lemma}\label{lemma: centralizers give closed subroot systems}
Assume we are in the Setup \ref{setup: finite and affine root systems}. Let $G_{\Psi^{\bullet}}$ be the corresponding group and $T_{\Psi^{\bullet}}$ its maximal torus.
Let $y \in T_{\Psi^{\bullet}}$ be any closed point. Then its connected centralizer $\Cent^{\circ}_{G_{\Psi^{\bullet}}}(y)$ is a regular subgroup of $G_{\Psi^{\bullet}}$ with associated closed root subsystem 
\begin{equation*}
   \Psi^{\bullet}_{\Cent^{\circ}(y)} = \{ \beta \in \Psi^{\bullet} \mid \beta(y) = 1 \} \subseteq \Psi^{\bullet}.
\end{equation*}
\end{lemma}
\begin{proof}

Since we are working in characteristic zero and $\Cent^{\circ}(y)$ is connected, it is sufficient to describe its Lie algebra. This is clearly a regular subalgebra -- it is stable under the action of the Cartan subalgebra corresponding to the maximal torus $T$ containing $y$. Hence it comes from a closed root subsystem by Proposition \ref{proposition: closed root subsystems and regular subalgebras}. The subsystem is given by $\Psi^{\bullet}_{\Cent^{\circ}(y)}$, as can be checked root-by-root: the adjoint action of $y$ fixes the root space $U_{\beta}$ labeled by $\beta$ if and only if $\beta(y)=1$.
\end{proof}

\begin{remark}
It is also easy to check that $\Psi^{\bullet}_{\Cent^{\circ}(y)}$ is a closed root subsystem by hand. Indeed, assume $\alpha, \beta \in \Psi^{\bullet}_{\Cent^{\circ}(y)}$. Then $\alpha(y) = \beta(y) =1$.
%
To see that $\Psi^{\bullet}_{\Cent^{\circ}(y)}$ is a root subsystem, note that $s_{\alpha}(\beta) = \beta {-\langle \beta, \alpha^{\vee} \rangle}\cdot \alpha$. Hence also $s_{\alpha}(\beta)(y) = \beta(y) \cdot \alpha(y)^{-\langle \beta, \alpha^{\vee} \rangle} = 1$, showing $s_{\alpha}(\beta) \in \Psi^{\bullet}_{\Cent^{\circ}(y)}$. 
%
To see that $\Psi^{\bullet}_{\Cent^{\circ}(y)}$ is closed, note that $(\alpha + \beta)(y) = \alpha(y) \cdot \beta(y) = 1$, showing $\alpha + \beta \in \Psi^{\bullet}_{\Cent^{\circ}(y)}$.
\end{remark}




We review the structure of these centralizers in more detail in the two cases from Setup \ref{setup: finite and affine root systems} separately in \S \ref{section: finite case and pseudo-levi subgroups} and \S \ref{section: centralizers in the affine case} below.





\subsubsection{Finite case and pseudo-Levi subgroups.}\label{section: finite case and pseudo-levi subgroups}
We now specialize to the finite case.
%
Let $G$ be a reductive group over an algebraically closed field $k$. Given any semisimple element $x \in G$, we wish to recall the structure of its centralizer $\Cent_{G}(x)$. 
%
Let $T \leq G$ be a maximal torus. Clearly, $\Cent_G(x)$ depends only on the conjugacy class of $x$. Conjugating inside $G$, we may thus without loss of generality assume $x \in T$.


\begin{proposition}\label{proposition: pseudo-levi subgroups}
Let $G$ be a connected reductive group over and algebraically closed field $k$. Given any $x \in T$, we have the following.
\begin{enumerate}
    \item[(i)] The centralizer $\Cent_{G}(x) \leq G$ of $x$ in $G$ is a reductive group.
    \item[(ii)] It is generated by $T$, the root subgroups $U_{\beta}$ for those roots $\beta \in \Phi^{\bullet}$ such that $\beta(x)=1$, and those Weyl group representatives $n_w$ that commute with $x$.
    \item[(iii)] The identity component $\Cent^{\circ}_{G}(x)$ is a connected reductive group with root system determined by $(X^{\bullet}, \Phi_{\Cent^{\circ}(x)})$, where $X_{\bullet} = X_{\bullet}(T)$ and $\Phi^{\bullet}_{\Cent^{\circ}(x)} = \{ \beta \in \Phi^{\bullet}_G \mid \beta(x)=1\}$. Its Weyl group $W_x \leq W$ is the subgroup generated by the reflections $s_{\beta}$, $\beta \in \Phi^{\bullet}_{L_x}$.   
    \item[(iv)]  If $G$ is moreover simply connected, $\Cent_{G}(x)$ is connected.     
\end{enumerate}
\end{proposition}
\begin{proof}
See \cite[\S 2]{Hum95} for an excellent overview. 
Statements (i), (ii), (iii) are \cite[\S 2.2. Theorem]{Hum95}; (iv) is \cite[\S 2.11. Theorem]{Hum95}.
%
Also see \cite[\S 4, \S 5]{Ste68}, \cite[\S 4, in particular \S 4.1 and \S 4.2]{SS70}, \cite[\S 6]{MS03}.
\end{proof}

\begin{notation}
The connected centralizers $\Cent^{\circ}(x) \leq G$ for varying $x\in T$ are called \emph{pseudo-Levi subgroups} of $G$. These are regular subgroups of $G$.
\end{notation}


\begin{remark}
Let $G$ be a reductive group and $x \in T$.
\begin{enumerate}
\item[(1)] Assume $G$ is of type A. Then any pseudo-Levi $\Cent^{\circ}(x)$ is in fact a Levi subgroup of $G$. Its Weyl group $W_x$ is generated by a subset of simple reflections.
%
\item[(2)] In other types, $\Cent^{\circ}(x)$ might not be a Levi subgroup of $G$. See \cite[Table 1]{RV19}, \cite[\S 12.1]{Kan01}. for a complete list of pseudo-Levi subgroups that appear in a given $G$.
%
This classification was done in \cite{BD49}, by describing all (closed) subroot systems of $\Phi^{\bullet}_G$.
%
\item[(3)] If $G$ is not simply connected, the centralizer $\Cent(x)$ might not be connected. For example \cite[4.3.Examples]{SS70}, the centralizer of $(i, -i)$ in $PGL_2$ is given by two discrete points corresponding to monomial matrices.
\end{enumerate}  
\end{remark}


\subsubsection{Centralizers in the affine case.}\label{section: centralizers in the affine case}
In the affine case, the above discussion immediately gives the following.

\begin{lemma}\label{proposition: affine centralizer as regular subgroup}
Let $\lo G \rtimes \Gm^{\rot}$ be the loop group of a reductive group $G$ extended by loop rotation. Take $(x, \zeta) \in T \times \Gm^{\rot}$ and denote $\ell =\ord(\zeta) \geq 0$. 
The identity component $\Cent^{\circ}_{\lo G \rtimes \Gm^{\rot}}(x, \zeta)$ is a regular subgroup of $\lo G$ associated to the closed root subsystem $\Phi^{\bullet}_{\Cent^{\circ}(x, \zeta)} = \{ \beta + m\hbar \in \Phi^{\bullet}_{G, \aff} \mid \beta(x) \cdot \zeta^m = 1 \}$.
\end{lemma}
\begin{proof}
Firstly, note that 
\begin{equation*}
    \Cent^{\circ}_{\lo G \rtimes \Gm^{\rot}}(x, \zeta) = \Cent^{\circ}_{\lo G}(x, \zeta) \rtimes \Gm^{\rot}.
\end{equation*}
because for any $\zeta' \in \Gm^{\rot}$, the elements $(x, \zeta)$ and $(1, \zeta')$ lie in the same torus $T \times \Gm^{\rot}$ and thus commute. Hence we can restrict to the centralizer inside $\lo G$. 

Since we consider the connected centralizer, it is enough to describe its affine root spaces.
The conjugation by $(x, \zeta)$ acts on such root space $U_{\beta + m\hbar}$ by the scalar $\beta(x) \cdot \zeta^m$ as in Lemma \ref{lemma: weights of extended torus on affine root groups}.
The fixed points are the whole $U_{\beta + m\hbar}$ if $\beta(x) \cdot \zeta^m =1$ and a single point else. The statement follows.   
\end{proof}
%
We will return to Lemma \ref{proposition: affine centralizer as regular subgroup} and give a more detailed description of $\Cent^{\circ}_{\lo G}(x, \zeta)$ and it root system later in Lemma \ref{lemma: root spaces of J} and Lemmata.

\begin{remark}
Alternatively, \cite{RV19} classified all closed root subsystems of $\Phi^{\bullet}_{G, \aff}$, see in particular \cite[Proposition 4.1.1 and Theorem 4.2.1]{RV19}; also see \cite[Theorem 5.3]{FRT08}.
The computations above can be also bootstrapped from \cite[Proposition 3.2.1 with notation of \S 3.1]{RV19} 
\end{remark}


\begin{remark}
The closely related classification problem for all (not necessarily closed) root subsystems of $\Phi^{\bullet}_{G, \aff}$ was resolved much earlier \cite{Dye90, DL11a, DL11b}. The answer here is similar, but slightly more complicated \cite[Theorem 5.1 and \S 5.5]{Dye90}, \cite[Theorem 3]{DL11a}.
\end{remark}




\subsubsection{Stabilizers in Weyl groups as reflection subgroups.}\label{section: stabilizers in weyl groups as reflection subgroups}
Let $W_{\Psi^{\bullet}_1}$ denote the Weyl group of the ambient root system $\Psi^{\bullet}_1$. It has the structure of a Coxeter group, whose reflections are labeled by the elements of $\Psi^{\bullet}_1$.
\begin{notation}
A subgroup $W'$ of $W_{\Psi^{\bullet}_1}$ is called a \emph{reflection subgroup} if it is generated by a subset of reflections $S' \subseteq (s_{\beta} \mid \beta \in \Psi^{\bullet}_1)$.    
\end{notation}


\begin{remark}
Reflection subgroups of $W_{\Psi^{\bullet}_1}$ correspond to (not necessarily closed) root subsystems of ${\Psi^{\bullet}_1}$. Indeed, to such $W'$ we can associate the subset $\{\beta \in \Psi^{\bullet}_1 \mid s_{\beta} \in W'\}$, which is a root subsystem of $\Psi^{\bullet}_1$. Vice versa, to any root subsystem we can associate the reflection subgroup generated by reflections in its roots. This defines a bijection between root subsystems $\Psi^{\bullet}_2$ of $\Psi^{\bullet}_1$ and reflection subgroups $W_{\Psi^{\bullet}_2}$ of its Weyl group $W_{\Psi^{\bullet}_1}$. See \cite[\S 1.2]{DL11a}, \cite[\S 2]{FRT08}). 
\end{remark}

Assume we are in the Setup \ref{setup: finite and affine root systems}.
For simplicity, assume that $G=G^{\sc}$ is simply connected.
The Weyl group $W_{\Psi^{\bullet}_1}$ naturally acts on the maximal torus $T_{\Psi^{\bullet}}$ of the corresponding group -- in the two respective cases we have $W_G \acts T$ and $W_{G, \aff} \acts T \times \Gm^{\rot}$. The following result is classical.

\begin{proposition}\label{proposition: steinberg fixed point theorem}
Assume we are in the Setup \ref{setup: finite and affine root systems} and $G=G^{\sc}$. Let $ z \in T_{\Psi^{\bullet}}$ be any closed point of the maximal torus. Then the stabilizer of $z$ under the action $W_{\Psi^{\bullet}_1} \acts T_{\Psi^{\bullet}_1}$ is a reflection subgroup of $W_{\Psi^{\bullet}_1}$.
\end{proposition}
\begin{proof}
This is due to Steinberg; variants of this result are called \emph{Steinberg fixed point theorems}. In the case of finite Weyl group $W$, this is basically \cite[Theorem 1.5]{Ste64}. The statement in terms of tori, including the generalization to the affine case of $W_{\aff} = \Z\Phi^{\bullet}_G \rtimes W$, follows from \cite[\S 3, \S 4, \S 5]{Ste68} -- see in particular \cite[\S 4.2]{Ste68} for the relevant statement on fixed-points on real topological tori and \cite[\S 5.1 and \S 5.3 Theorem]{Ste68} for the translation to algebraic tori. The simply-connectedness is needed so that $X_{\bullet} = \Z\Phi_{\bullet}$.

For alternative approaches to the finite case, see \cite[chap. V]{Bou68} or \cite[Theorem 1.1]{Leh04}. For the general case, see \cite[Theorem 2.2]{PS19}, \cite[Introduction]{Pue23}. Further see \cite[\S 11, \S 12]{Kan01}.
\end{proof}

\begin{remark}
Recall the we assume $G=G^{\sc}$.
In Proposition \ref{proposition: steinberg fixed point theorem}, the stabilizer is the reflection subgroup corresponding to the closed root subsystem $\Phi^{\bullet}_{\Cent^{\circ}(z)}$, namely
\begin{equation}\label{equation: affine stabilizer as reflection group}
    \Stab_{W_{\Psi^{\bullet}_1}}(z) = W_{\Phi^{\bullet}_{\Cent^{\circ}(z)}}.
\end{equation}

Indeed, knowing Proposition \ref{proposition: steinberg fixed point theorem}, it suffices to describe the subset of reflections which lie inside the stabilizer. (Alternatively, we may lift elements of $W_{\Psi^{\bullet}_1}$ to $G_{\Psi^{\bullet}_1}$ and use the description of the stabilizer there.) We discuss this concretely in Lemmata \ref{lemma: stabilizer in finite weyl group}, \ref{lemma: affine stabilizer in weyl group}.
%
However, note that \eqref{equation: affine stabilizer as reflection group} may fail without the simply connected assumption; see Proposition \ref{proposition: pseudo-levi subgroups}, (ii). 
\end{remark}